\section{Additive source deepening after the first two goal pairs}
\label{r28:app:source-supplement}
The initial survey remains \textbf{30 formal reading records at 38 distinct URLs}. The later source-only supplement adds three attributed records: further selected passages of the already known Del Vecchio--Fr\"ohlich--Pizzo paper, an ETH institutional publication record, and a targeted topology reading of the already known Yarotsky paper. These are not three new independent studies. Together the current ledgers contain \textbf{33 formal reading records at 41 distinct recorded URLs}, plus five screening or reopen records; the URL union including those screens is 46. The intermediate two-record supplement had totals 32/39 and a screened URL union of 44; its creation checkpoint remains preserved. The one unread lead stays outside the reviewed count. No government-primary document was newly read in this supplement, and it adds zero research investigations.

The technical text is \href{https://arxiv.org/abs/2108.13907}{arXiv:2108.13907v1}, a 54-page preprint. The \href{https://www.research-collection.ethz.ch/items/142ea4e3-1538-4f05-8a97-5200a860d558}{ETH institutional record} verifies publication in \emph{Journal of Functional Analysis} 284(1), 109734 (2023), \href{https://doi.org/10.1016/j.jfa.2022.109734}{doi:10.1016/j.jfa.2022.109734}. Final journal text was not compared with v1. The recorded technical passages are Sections 1.1--1.2 and Remark 1.1; the Section 2 opening; definitions (2.29)--(2.32); Lemma 2.8's statement and initial proof; selected Section 2.4 gap estimates through Corollary 2.12; the Section 3 induction introduction and selected Theorem 3.1 inequalities; Theorem 3.3 and its closing proof; Lemma 4.1's statement; and Lemma 4.3 with the adjoining inverse-series justification. The long induction and appendices were not audited end to end.

The Feynman-method deepening initially compared an application to the original finite-range homogeneous Hamiltonian $H_0+\Phi$ with organizing AG3's generated interaction inside the cited paper's evolving-reference induction. Its later archive correction identifies an omitted established result: the current Round21 I1 admission, revision 2, already supplies the complete 24-link/21-face dictionary, qualitative original-model stability and the specified positive-orthant full-link-observable GNS construction. With source constants left symbolic,
\begin{equation}
 \tau_* =\frac{\min\{c_1(S),1/[2c_2(S)]\}}7>0,\qquad
 0<|\tau|<\tau_*\ \Longrightarrow\ \Delta\ge\alpha/16.
\end{equation}
I1 certifies no selected nonzero numerical $\tau$. Its homogeneous \emph{omitted} coupling on the selected-strip reference is not a uniform coefficient on every plaquette, a summable-model representation, or a continuum construction. The specified orthant GNS conclusion is not independence of every boundary sequence.

Consequently another theorem application giving only unspecified-small-coupling existence would duplicate the admitted qualitative outcome. A distinct new target must supply a numerical interval, complete generated-interaction iteration, or an additional representation or boundary statement beyond I1. The archive comparison uses the current repository admission and source dictionary; it neither rereads nor reconstructs the external Yarotsky proof, and adds no primary reading record or investigation. The unchanged deeper-source note and this additive correction are separately bound in \repo{research/round28/experts/source-supplement.json}.

A different published theorem would need its own hypothesis matching while reusing I1's established site, star, domain and normalization dictionary. Infinite onsite dimension is permitted by the Del Vecchio paper's stated assumptions; its initial range condition is not confined to the expository nearest-neighbor presentation. Theorem-specific boundary and gauge-observable restrictions must remain explicit, and a constrained physical space cannot silently become a tensor product. Re-entry with AG3's generated interaction additionally requires a proved passage from arbitrary indexed supports to rectangular energy-weighted data, changed local gaps, retained scalars and diagonals, and the complete induction. These are inferred transfer obligations, not defects asserted in either external theorem or missing work already discharged by I1.

The third supplemental reading, record R28-FGNS1, inspects Yarotsky's 17-page \href{https://arxiv.org/pdf/math-ph/0411042}{math-ph/0411042v1 PDF} and corresponding \href{https://arxiv.org/html/math-ph/0411042v1}{HTML}. Its depth is the definitions and Theorems 1--3 on printed pages 2--4, creation operators on pages 6--7, selected expansion and state-limit passages on pages 8--9, and the ending generator/core argument on pages 9--10. References on page 17 were bibliographic pointers only. The long proof and cited predecessor papers were not independently reconstructed; an older I1 download hash is not claimed as a hash of these newly inspected browser bytes.

The source states convergence of finite-ground expectations for each observable in the full local algebra $\bigcup_F B(\mathcal H_F)$, and resolvent matrix-element convergence tested on local excitation vectors. Its generator discussion uses local bounded commutators and creation vectors with $u_I\in\mathcal H'_I\cap D(H_{I,0})$, followed by density and essential self-adjointness statements. These are not trace-norm convergence of local density matrices, strong resolvent convergence in the earlier summable representation, or a declaration that every bounded-local cyclic vector lies in the generator domain. No explicit local-normality or gauge-average compatibility claim was located in the inspected passages; this limited reading does not prove that such properties fail or cannot follow from the full source theory.

For a later homogeneous physical-GNS identification, the source note therefore distinguishes pointwise state limits and a generator core from the regularity needed to pass local weak-operator Haar averages through the actual representation. It calls for local normality or another explicit continuity-and-averaging argument, continuous gauge implementers on all endpoint groups, the full bounded local invariant algebra, and a justified reducing spectral restriction. I1 already supplies the specified homogeneous full-algebra state and generator; J2's cyclic/invariant identification belongs to its dyadic summable representation and does not automatically transfer. These are prospective transfer obligations only. The reading constructs no vacuum-reset bound, local tightness estimate, gauge-cyclic equality or new gap.

The AG3 passive comparison and spent $9/4\to2\to3/2$ margins remain unchanged. AG2/AG3's complete-remainder and one-step results do not evaluate I1's source constants; their numerical extracted-reference gaps are not new numerical full-Hamiltonian gap theorems. The supplement supplies no new estimate, fixture, threshold, theorem application or continuum premise. It informed later planning without selecting a fourth or fifth goal at the time of the review. The source and reading-depth records, including the correction, are included in the manuscript source manifest.

The later admitted AJ1 investigation, Section~\ref{r28:sec:aj1}, now supplies the previously prospective local-energy and normality argument, compatible finite-endpoint averaging, physical cyclic/fixed equality and closed reducing restriction for the actual I1 state. Its post-review clarification uses local normality and finite spectator columns to justify represented product-vector density; it does not assume operator-norm density of the algebraic tensor product. This mathematical admission is distinct from the earlier source-only note. AJ1's primary retrieval records and restricted proof-reading claims are bound through its gate; they do not silently change the 33-record survey ledger. AJ2, Section~\ref{r28:sec:aj2}, then proves a nonzero homogeneous Wilson fluctuation using this normality and the actual multiplier's Haar-null levels. The coupling interval remains symbolic, and neither investigation evaluates a mass or the source stability constants.
